% Information Acquisition Channel

The baseline results establish that active funds reduce positions in stocks they lend. We now investigate the mechanism: do funds trade on information acquired through the lending process? We exploit mandatory short position disclosure events---as described in Section~\ref{subsec:identification-disclosure}---as an external validation of the information channel. The key identification test is whether active lenders reduce positions specifically in the month \textit{before} a large short position is publicly disclosed, consistent with acting on private information about the short seller's position.

\subsection{Mandatory Short Position Disclosure}

Table~\ref{table_disclosure} presents the results, estimated on the disclosure sample (observations within $\pm 3$ months of a first-time large short position disclosure for stocks that are ever publicly disclosed). The dependent variable is $PosChange_{i,j,t}$. $OnLoan_{i,j,t-1}$ is standardized lending activity; $\mathit{Disclosure}_{i,t+1}$ equals one if stock $i$ has a \textit{new} public short position disclosure in month $t+1$ (defined as having no prior disclosure in the preceding six months). All specifications use active funds and double-cluster standard errors by fund and time.

Column~(1) reports the baseline lending effect for this restricted sample: $OnLoan_{i,j,t-1} = -0.2837$, $t = -2.85$, confirming the baseline result holds in the disclosure sample. Column~(2) adds the disclosure indicator, which is positive but insignificant ($0.3633$, $t = 1.25$), consistent with an absence of general reaction to impending disclosures among both lenders and non-lenders alike (this effect is absorbed by the stock$\times$time fixed effects in the saturated specification).

Column~(3) introduces the key interaction: $OnLoan_{i,j,t-1} \times \mathit{Disclosure}_{i,t+1}$. The interaction coefficient is $-0.8027$ and statistically significant at the 5\% level ($t = -2.68$). Active funds reduce their positions substantially more---by an additional 0.80 percentage points per standard deviation of lending---when their lending activity precedes a new disclosure event by one month. Column~(4) adds the full saturated three-way fixed effects; the interaction coefficient strengthens to $-0.8622$ ($t = -2.22$), allaying concerns that the result is driven by unobserved stock-time or fund-time factors.

Column~(8) provides the critical placebo test. Estimating the identical specification for passive funds only, the coefficient on $OnLoan_{i,j,t-1}$ is economically small and insignificant ($0.0208$, $t = 0.23$), and the interaction with $\mathit{Disclosure}_{i,t+1}$ is precisely zero ($0.0258$, $t = 0.28$). Passive funds show no trading reaction to lending signals around disclosure events. This placebo result strongly supports the information acquisition interpretation: only funds with discretion to trade on private information do so.

\subsection{Passive Fund Lending as Alternative Measure}

A potential concern is that a fund's own lending activity is endogenous to its trading decisions---a fund might selectively recall shares or reduce positions for reasons correlated with lending. To address this, we replace $OnLoan_{i,j,t-1}$ with $\mathit{OnLoan}_{i,f\text{-passive},t-1}$, the aggregate lending of stock $i$ by passive funds in the same family (Table~\ref{table_passive_iv}). Since passive funds track indices and cannot trade on information, their lending is driven by portfolio composition and borrower demand---not strategic timing---providing an exogenous source of short selling demand signals available to the family.

The coefficient on passive fund lending is $-0.3406$ ($t = -3.20$) in column~(1), comparable to the baseline. The interaction with $\mathit{Disclosure}_{i,t+1}$ in column~(3) is $-1.1240$ and significant at the 1\% level ($t = -3.29$)---larger in magnitude than the baseline interaction ($-0.8027$), consistent with passive fund lending providing a cleaner, more exogenous measure of short selling demand. Column~(4) confirms robustness with saturated fixed effects ($-0.9156$, $t = -2.32$). These results indicate that the trading response to lending is not driven by endogenous selection of which stocks to lend.

\subsection{Timing Tests: Leads and Lags of Disclosure}

If active funds genuinely acquire information through lending and trade in anticipation of future short-selling pressure, only the \textit{lead} disclosure should predict trading. Contemporaneous and lagged disclosures should be uninformative once the information has become public.

Table~\ref{table_timing} tests this prediction by including interactions of $OnLoan_{i,j,t-1}$ with $\mathit{Disclosure}_{i,t+1}$ (lead), $\mathit{Disclosure}_{i,t}$ (contemporaneous), and $\mathit{Disclosure}_{i,t-1}$ (lag). Column~(3) shows that only the lead interaction is statistically significant ($-0.8224$, $t = -2.62$). The contemporaneous interaction is economically negligible ($0.0353$, $t = 0.08$) and the lagged interaction is small and insignificant ($-0.1157$, $t = -0.62$). Column~(4) with saturated fixed effects confirms this pattern: the lead interaction is $-0.8723$ ($t = -2.27$), while contemporaneous and lagged interactions remain insignificant.

This timing pattern---pre-announcement selling by lenders but not in response to past or contemporaneous disclosure---is precisely what the information acquisition channel predicts. Active lending funds reduce positions in the month before a large short position becomes publicly known, consistent with observing borrowing demand before it crosses the public disclosure threshold.
